For another proper primary \(\mathfrak b\), put \(d(v)=v_{\mathfrak a}(\mathfrak b)>0\), \(T(v)=v/d(v)\). G331 supplies a homeomorphism of sections and bounded continuous \(d,d^{-1}\). Transport degree-one homogeneous inputs by \(u_{\mathfrak b}(Tv)=u_{\mathfrak a}(v)/d(v)\). Rational profiles have the same ideals and denominators in both gauges, while tests satisfy \(v_{\mathfrak b}(J)(Tv)=v_{\mathfrak a}(J)(v)/d(v)\). The proposed measures have identical ideal moments, since the factors cancel:
\[
\mu_{\mathfrak b}(\boldsymbol u_{\mathfrak b})
=T_*\big[d\,\mu_{\mathfrak a}(\boldsymbol u_{\mathfrak a})\big]. \tag{GAUGE}
\]
Density proves this first on rational tuples. Uniform \(C\) and actual bounded decreasing \(D\) continuity extend it: transport preserves these convergences, and multiplication by fixed \(d\) and compact pushforward are weakly continuous. The factor is ONE \(d\), independently of dimension. The new mass is the old moment against \(v_{\mathfrak a}(\mathfrak b)\); for rational tuples it is \(e(I_1,\ldots,I_p,\mathfrak b)/\prod q_i\). The cocycle follows because successive weights multiply to \(v_{\mathfrak a}(\mathfrak c)\). On geometric atoms, \(\operatorname{ord}_Z(\mathfrak b)=d(\nu_Z)b_Z\) gives the same law directly.

Let \(S\) be the finite semilocal normalization, \(q\) its normal local centers above \(\mathfrak m\), \(f_q=[\kappa(q):\kappa(\mathfrak m)]\), and \(\rho_q\) restriction on the full boundary normalized by \(\mathfrak aS_q\). A regular conductor element kills \(S/R\). Any prime in its support has dimension below \(n\), since a chain above it extends strictly down to a minimal prime not containing that element. N1's exact-sequence top-form additivity therefore gives the same degree-\(n\) length form for \(R\) and \(S\). For each primary product \(P\), Artinian decomposition and composition series give
\[
\ell_R(S/PS)=\sum_q f_q\,\ell_{S_q}(S_q/PS_q).
\]
The residue degrees are field dimensions, including inseparable extensions. Comparing top forms yields
\[
e_R(I_1,\ldots,I_n)
=\sum_{\dim S_q=n}f_q\,e_{S_q}(I_1S_q,\ldots,I_nS_q).
\]
Testing on the dense ideal lattice, then extending by the two continuities, proves
\[
\mu_{R,\mathfrak a}(\boldsymbol u)
=\sum_{\dim S_q=n}f_q\,\rho_{q,*}
 \mu_{S_q,\mathfrak aS_q}(\rho_q^*\boldsymbol u). \tag{NORMAL}
\]
At lower-dimensional centers the degree-\(n\) contribution is ZERO; it is not their own lower-arity intrinsic measure. They and every support stay in the spaces and potential domain. There is no point deletion, nor an invented generic degree for a vertical component.

G331 describes the FULL normalization quotient: conductor-finite fibers are singletons; conductor-infinite fibers are exactly equal restrictions through \(S/\mathfrak c\to R/\mathfrak c\). It supplies positive \(C,D\) descent precisely for full-fiber-constant functions. Thus any compatible upstairs tuple in these classes gives actual downstairs inputs, and (NORMAL) computes their product by pushing along the full quotient, adding any masses on identified fibers. Both strata, nonminimal supports and mixed-zero patterns remain. This is a normalization/conductor rule, not arbitrary finite-map or envelope functoriality.

For an excellent reduced curve \(n=1,p=0\), normal centers are DVRs. A semivaluation there is a positive multiple of its order (nonzero elements are a uniformizer power times a unit); its only allowable support is zero. A regular conductor element has finite value on every branch, so the entire downstairs boundary is the finite branch set
\[
\nu_q=\operatorname{ord}_q/b_q,\qquad b_q=\operatorname{ord}_q(\mathfrak a)>0.
\]
The independently geometric zero-input measure is
\[
\mu_{R,\mathfrak a}()=\sum_q f_qb_q\,\delta_{\nu_q},\qquad
\mu_{R,\mathfrak a}(X)=e_R(\mathfrak a). \tag{CURVE}
\]
Here branches mean normal LOCAL centers, including multiple branches of one irreducible component.

Its actual Cartier comparison can also be proved directly on every curve primary model. Since \(IS\) is invertible, \(\operatorname{Spec}S\to Y\) exists. It is finite by its closed graph in \(\operatorname{Spec}S\times_RY\), and surjective because its finite image contains every component generic point of reduced \(Y\). It is the finite normalization. For effective exceptional Cartier \(E=(g)\) at a closed point, put \(B=\mathcal O_{Y,y}\), \(\widetilde B\) its finite semilocal normalization, \(C=\widetilde B/B\) of finite length. The snake sequence for multiplication by regular \(g\) gives
\[
0\to C[g]\to B/gB\to\widetilde B/g\widetilde B\to C/gC\to0.
\]
On finite-length \(C\), the kernel and quotient have equal lengths. Therefore the local Cartier degree, retaining residue weights, equals the sum of DVR orders with their residue degrees. Summing over \(y\), extending by effective differences and real coefficients, gives
\[
Q_Y(F)=\sum_q f_q\operatorname{ord}_q(F)
=\int_X H_F\,d\mu_{R,\mathfrak a}().
\]
This is an actual supported Cartier degree, agrees with n3--n4, and is independent of models. DVR lengths give \(\int v(J)\,d\mu=e_R(J)\). Zero-arity locality and continuity are constant identities.

Residue weights cannot be replaced by generic degree. Let \(k\subset l\) be ANY finite extension of degree \(d>1\), possibly inseparable, and \(R=k+t\,l[[t]]\subset S=l[[t]]\). If \(1,\alpha_2,\ldots,\alpha_d\) is a \(k\)-basis, \(R=k[[t]]\oplus\bigoplus_{i=2}^d t\,k[[t]]\alpha_i\) is a finite algebra over the excellent DVR \(k[[t]]\), hence excellent Noetherian. It is local, with \(\mathfrak m=tS\), residue \(k\), fraction field \(l((t))\), finite normalization \(S\), conductor \(tS\), generic degree one and residue degree \(d\). For \(J=t^rS\), a residue basis gives
\[
\ell_R(R/J^N)=1+d(rN-1),\qquad e_R(J)=dr.
\]
For gauge \(\mathfrak m\), (CURVE) is \(d\delta_{\operatorname{ord}_t}\), with the correct moment \(dr\), regardless of separability.

If \(n=0\), the reduced local ring is a field, its centered boundary is empty, and its only measure is zero; no negative-arity operator or degree-zero multiplicity mass is asserted. For \(n\ge2\), a zero input (represented by a unit) annihilates the product; a unit test has zero moment. For \(p=0\) there is no input zero to annihilate the reference measure. The tuple \(-t_i v(\mathfrak a)\) has mass \((\prod_i t_i)e_R(\mathfrak a)\), with empty product one. A zero ideal has value \(+\infty\), gives an unbounded bottom profile, and is outside the bounded product domain and outside primary gauges; it must not be multiplied formally by scalar zero. Define a zero scalar input as the actual zero function. Unit or zero ideals have no level-one section.

The reduction convention is necessary. \(k[[t]][\varepsilon]/(\varepsilon^2)\) and \(k[[t]]\) have the same semivaluation boundary, but parameter-power lengths are \(2s\) and \(s\). Nonreduced generic lengths require extra data. We claim reduced multiplicities only, while the genuine nilpotent thicknesses of Cartier divisors on reduced models were fully retained in n3.
